\begin{problem}[第34届物理竞赛复赛][彗星轨道]{73}
    星体 $P$（行星或彗星）绕太阳运动的轨迹为圆锥曲线
    \begin{equation}
        r = \frac{k}{1 + \varepsilon \cos \theta}
    \end{equation}
    式中，$r$ 是 $P$ 到太阳 $S$ 的距离，$\theta$ 是矢径 $SP$ 相对于极轴 $SA$ 的夹角（以逆时针方向为正），$k = L^{2}/(GMm^{2})$，$L$ 是 $P$ 相对于太阳的角动量，$G = 6.67 \times 10^{-11}\,\mathrm{m^{3}\cdot kg^{-1}\cdot s^{-2}}$ 为引力常量，$M \approx 1.99 \times 10^{30}\,\mathrm{kg}$ 为太阳的质量，$\varepsilon = \sqrt{1 + \frac{2EL^{2}}{G^{2}M^{2}m^{3}}}$ 为偏心率，$m$ 和 $E$ 分别为 $P$ 的质量和机械能。假设有一颗彗星绕太阳运动的轨道为抛物线，地球绕太阳运动的轨道可近似为圆，两轨道相交于 $C$、$D$ 两点，如图所示。已知地球轨道半径 $R_{E} \approx 1.49 \times 10^{11}\,\mathrm{m}$，彗星轨道近日点 $A$ 到太阳的距离为地球轨道半径的三分之一，不考虑地球和彗星之间的相互影响。求彗星
    \begin{subproblem}
        先后两次穿过地球轨道所用的时间；
    \end{subproblem}
    \begin{subproblem}
        经过 $C$、$D$ 两点时速度的大小。
    \end{subproblem}
    已知积分公式 $\displaystyle \int \frac{x\,\mathrm{d}x}{\sqrt{x+a}} = \frac{2}{3}(x+a)^{3/2} - 2a(x+a)^{1/2} + C$，式中 $C$ 是任意常数。
\end{problem}